% =====================================================================
%  Chapitre 3 — Réalisation
% =====================================================================
\chapter{Réalisation}
\label{chap:realisation}

\section{Environnement et technologies}

Les versions ci-dessous sont celles qui figurent dans \code{backend/requirements.txt} et
\code{frontend/package.json}. Lorsqu'une version n'est pas épinglée dans un fichier du
projet, la mention correspondante l'indique.

\begin{longtable}{@{}L{4.2cm} L{2.9cm} L{7.1cm}@{}}
\caption{Technologies employées et versions réellement déclarées}
\label{tab:technos}\\
\toprule
\textbf{Technologie} & \textbf{Version} & \textbf{Rôle dans le projet} \\
\midrule
\endfirsthead
\multicolumn{3}{c}{\small\itshape suite du tableau~\ref{tab:technos}}\\
\toprule
\textbf{Technologie} & \textbf{Version} & \textbf{Rôle} \\
\midrule
\endhead
\bottomrule
\endfoot

\multicolumn{3}{@{}l}{\textit{\textbf{Backend}}}\\
Python      & 3.12.7 & Langage du serveur (version de l'interpréteur utilisé) \\
FastAPI     & 0.115.6 & Framework web, validation des entrées, OpenAPI \\
Uvicorn     & 0.34.0 & Serveur ASGI \\
SQLAlchemy  & 2.0.36 & ORM, style déclaratif typé \\
Alembic     & 1.14.0 & Migrations de schéma versionnées \\
psycopg2-binary & 2.9.10 & Pilote PostgreSQL \\
Pydantic    & 2.10.4 & Schémas d'entrée et de sortie \\
pydantic-settings & 2.7.0 & Configuration par variables d'environnement \\
python-jose & 3.3.0 & Signature et vérification des jetons JWT \\
passlib     & 1.7.4 & Interface de hachage \\
bcrypt      & 4.0.1 & Algorithme de hachage des mots de passe \\
email-validator & 2.2.0 & Validation des adresses électroniques \\
python-multipart & 0.0.20 & Formulaire du flux OAuth2 \\
pytest      & 8.3.4 & Exécution des tests \\
httpx       & 0.28.1 & Client HTTP de test \\
\midrule
\multicolumn{3}{@{}l}{\textit{\textbf{Frontend}}}\\
Node.js     & 24.12.0 & Environnement d'exécution (version utilisée) \\
npm         & 11.6.2 & Gestionnaire de paquets \\
React       & 18.3 & Bibliothèque d'interface \\
TypeScript  & 5.7 & Typage statique \\
Vite        & 6.0 & Serveur de développement et empaqueteur \\
React Router & 6.28 & Routage côté client \\
TanStack Query & 5.62 & Cache et synchronisation des données serveur \\
Recharts    & 2.13 & Graphiques \\
axios       & 1.7 & Client HTTP \\
\midrule
\multicolumn{3}{@{}l}{\textit{\textbf{Infrastructure}}}\\
PostgreSQL  & 16-alpine (Docker) & Base déclarée dans \code{docker-compose.yml} \\
PostgreSQL  & 12.4 (local) & Instance effectivement utilisée en développement \\
Docker Compose & non versionné & Orchestration des trois services \\
\end{longtable}

\begin{remarque}[Lecture des versions du frontend]
Les dépendances JavaScript sont déclarées avec l'opérateur \code{\textasciicircum{}},
qui autorise les montées de version mineures. Le tableau retient la version minimale
déclarée. La version exactement installée figure dans \code{package-lock.json}.
\end{remarque}

\section{Organisation du code}

\begin{figure}[H]
\centering
\begin{minipage}{0.92\textwidth}
\small\ttfamily\raggedright
\textcolor{medicaDeep}{backend/app/}\\
\hspace*{1em}core/ \hfill \textcolor{medicaMuted}{\rmfamily config, security (JWT, bcrypt), rbac (catalogue)}\\
\hspace*{1em}db/ \hfill \textcolor{medicaMuted}{\rmfamily session, base déclarative, seed}\\
\hspace*{1em}models/ \hfill \textcolor{medicaMuted}{\rmfamily 34 tables réparties en 7 modules}\\
\hspace*{1em}schemas/ \hfill \textcolor{medicaMuted}{\rmfamily schémas Pydantic d'entrée et de sortie}\\
\hspace*{1em}services/ \hfill \textcolor{medicaMuted}{\rmfamily 7 services métier, sans dépendance HTTP}\\
\hspace*{1em}api/deps.py \hfill \textcolor{medicaMuted}{\rmfamily session, utilisateur, permission, périmètre}\\
\hspace*{1em}api/routers/ \hfill \textcolor{medicaMuted}{\rmfamily 10 routers, 99 endpoints}\\
\hspace*{1em}main.py \hfill \textcolor{medicaMuted}{\rmfamily application, CORS, montage des routers}\\[4pt]
\textcolor{medicaDeep}{backend/alembic/versions/} \hfill \textcolor{medicaMuted}{\rmfamily 3 migrations}\\
\textcolor{medicaDeep}{backend/tests/} \hfill \textcolor{medicaMuted}{\rmfamily 7 fichiers, 112 tests}\\[6pt]
\textcolor{medicaDeep}{frontend/src/}\\
\hspace*{1em}lib/ \hfill \textcolor{medicaMuted}{\rmfamily client axios, contexte d'authentification, hooks}\\
\hspace*{1em}components/ \hfill \textcolor{medicaMuted}{\rmfamily mise en page, graphiques, panneaux, icônes}\\
\hspace*{1em}pages/ \hfill \textcolor{medicaMuted}{\rmfamily un composant par écran}\\
\hspace*{1em}styles.css \hfill \textcolor{medicaMuted}{\rmfamily charte graphique, variables CSS}
\end{minipage}
\caption{Arborescence des répertoires significatifs}
\label{fig:arborescence}
\end{figure}

La règle d'organisation du backend est simple : un router ne calcule rien. Il vérifie
les droits, charge les entités, délègue à un service, journalise et renvoie. Toute
formule est dans \code{services/}.

% =====================================================================
\section{L'intelligence de la plateforme : un moteur de règles, pas un modèle}
\label{sec:pas-dia}

Cette section est la plus importante du rapport en matière d'honnêteté technique. Le
produit s'appelle \emph{AOI Transformation}, une table de la base s'appelle
\code{ai\_recommendation}, et la plateforme recommande des types de solutions
d'intelligence artificielle. Rien de tout cela n'implique qu'elle en contienne une.

\subsection{Le constat, établi sur le code}

Une recherche exhaustive a été menée sur l'ensemble du dépôt, portant sur les
bibliothèques d'apprentissage automatique (\code{scikit-learn}, \code{tensorflow},
\code{torch}, \code{transformers}, \code{keras}, \code{xgboost}, \code{pandas},
\code{numpy}), sur les clients de fournisseurs de modèles (\code{openai},
\code{anthropic}, \code{mistralai}, \code{cohere}, \code{huggingface}), ainsi que sur
tout appel réseau sortant du backend. Le résultat est sans ambiguïté :

\begin{itemize}
  \item aucune de ces bibliothèques n'apparaît dans \code{requirements.txt} ni dans
    \code{package.json} ;
  \item aucun import correspondant n'existe dans le code ;
  \item aucun fichier de modèle entraîné (\code{.pkl}, \code{.h5}, \code{.onnx},
    \code{.pt}) n'est présent dans le dépôt ;
  \item aucune clé d'API de fournisseur de modèle n'est attendue par la configuration ;
  \item le backend n'émet aucune requête HTTP sortante. Le seul client HTTP installé,
    \code{httpx}, sert exclusivement aux tests.
\end{itemize}

Les deux seules occurrences du sigle « LLM » dans le code sont des chaînes de caractères
en français, décrivant le \emph{patron de solution} recommandé à l'utilisateur —
« Assistant LLM ancré sur base documentaire (RAG), sortie validée ». Ce sont des
recommandations adressées à un humain, pas des appels à un modèle.

\begin{remarque}[Formulation exacte à retenir]
La plateforme \emph{recommande} des solutions d'intelligence artificielle. Elle n'en
\emph{utilise} aucune. Son moteur de qualification est un système à base de règles
pondérées, entièrement déterministe, dont chaque sortie est reconstituable à la main.
\end{remarque}

\subsection{Moteur de règles et apprentissage automatique : la distinction}

La confusion entre les deux est fréquente, y compris dans la littérature commerciale.
Elle porte sur l'origine de la connaissance.

\begin{table}[H]
\centering
\caption{Ce qui sépare un moteur de règles d'un modèle d'apprentissage}
\label{tab:regles-vs-ml}
\small
\begin{tabular}{@{}L{3.3cm} L{5.3cm} L{5.3cm}@{}}
\toprule
& \textbf{Moteur de règles (ce projet)} & \textbf{Apprentissage automatique} \\
\midrule
Origine de la connaissance
& Écrite par un expert humain, lisible dans le code
& Inférée statistiquement depuis des données d'entraînement \\
Forme
& Prédicats booléens pondérés
& Paramètres numériques appris, souvent des millions \\
Besoin en données
& Aucun jeu d'entraînement
& Corpus annoté, représentatif et volumineux \\
Déterminisme
& Deux exécutions sur les mêmes faits donnent le même résultat
& Résultat dépendant de l'entraînement et du tirage \\
Explicabilité
& Native : la règle qui a décidé est nommée
& À reconstruire après coup (importance des variables, SHAP, LIME) \\
Évolution
& Modifier une règle, une ligne de code, relire la revue
& Réentraîner, réévaluer, surveiller la dérive \\
Validation réglementaire
& Le jeu de règles est un document versionnable et signable
& Exige un dossier de validation et des jeux de test figés \\
Limite principale
& Ne découvre rien qu'un expert n'ait prévu
& Capte des régularités qu'aucun expert n'aurait formulées \\
\bottomrule
\end{tabular}
\end{table}

Le moteur de ce projet relève de la première colonne. On peut le qualifier de système
expert au sens classique : une base de connaissances — 39~règles réparties sur 8~types
de solution — et un mécanisme d'inférence — l'agrégation pondérée et l'application des
règles bloquantes.

\subsection{Pourquoi ce choix est adapté ici}

Ce n'est pas un pis-aller. Quatre arguments le soutiennent.

\textbf{Il n'y a pas de données d'entraînement.} Entraîner un modèle à prédire
« quel type de solution convient à ce besoin » supposerait un historique de plusieurs
centaines de besoins qualifiés puis validés par leur résultat réel. Aucune organisation
ne dispose de ce corpus au moment où elle démarre sa démarche d'adoption. Le problème
est, par construction, un problème de démarrage à froid.

\textbf{La décision doit être opposable.} La recommandation est présentée en comité
d'investissement. Un expert du domaine doit pouvoir contester une règle précise. Avec
un modèle statistique, la discussion se déplace vers la qualité du jeu d'entraînement,
terrain sur lequel un comité métier ne peut pas arbitrer.

\textbf{Le contexte réglementaire pénalise l'opacité.} En périmètre GxP, une décision
non explicable est une décision à documenter lourdement. Un jeu de règles versionné est
un artefact de validation naturel.

\textbf{Le déterminisme rend le système testable.} Les règles sont couvertes par des
tests qui vérifient des sorties exactes. Un modèle probabiliste ne se teste pas ainsi ;
il se mesure sur des métriques agrégées.

\subsection{Ce que cela implique, et l'évolution envisageable}

La contrepartie est réelle : le moteur ne découvre rien. Il ne peut pas proposer un type
de solution auquel les règles n'ont pas pensé, ni apprendre des décisions passées. Si
tous les comités écartent systématiquement les recommandations « agentique », le moteur
continuera à les produire.

L'intégration d'un modèle de langage est envisageable et serait utile sur trois points
précis : pré-remplir les faits de qualification à partir d'une description libre du
besoin, rédiger un premier jet de fiche projet, et détecter des doublons sémantiques
entre opportunités formulées différemment. Ces trois usages sont décrits au
chapitre~\ref{chap:conclusion} comme \emph{perspectives}. Aucun n'est implémenté, et
aucun ne devrait remplacer la décision du moteur de règles : le modèle proposerait des
entrées, les règles continueraient de décider.

% =====================================================================
\section{Le moteur de qualification (M4)}

\subsection{Principe}

Le moteur compare une opportunité à huit types de solution. Il ne choisit pas « le plus
probable » ; il évalue chaque type indépendamment, puis retient celui dont le plus grand
poids de règles est satisfait.

Les huit types sont : IA classique (apprentissage supervisé, séries temporelles), vision
par ordinateur, traitement du langage, IA générative, IA agentique, décisionnel,
automatisation robotisée de processus, et refonte de processus. Ce dernier mérite d'être
souligné : c'est la réponse « pas d'IA », et le moteur la produit lorsque le processus
est instable ou les données absentes.

\subsection{Les faits d'entrée}

Quinze faits sont extraits de l'opportunité. Treize sont booléens
(\code{data\_available}, \code{data\_structured}, \code{ground\_truth\_available},
\code{unstructured\_text}, \code{image\_or\_video}, \code{content\_generation},
\code{multi\_step\_autonomy}, \code{rule\_based}, \code{repetitive\_task},
\code{reporting\_need}, \code{process\_unstable}, \code{regulated\_gxp},
\code{human\_in\_the\_loop}), deux sont numériques (\code{data\_quality} sur~5 et
\code{volume\_per\_month}). Les réponses libres du questionnaire peuvent surcharger un
fait, ce qui permet d'affiner sans modifier le schéma.

\subsection{Le calcul}

Soit $R_t$ l'ensemble des règles du type~$t$, $w_r$ le poids de la règle~$r$ et
$\sigma_r \in \{0,1\}$ l'indicateur de satisfaction. Le score brut du type vaut :

\begin{equation}
s_t = \frac{\displaystyle\sum_{r \in R_t} w_r \, \sigma_r}{\displaystyle\sum_{r \in R_t} w_r}
\label{eq:score-type}
\end{equation}

Une règle peut être déclarée \emph{bloquante}. Si une règle bloquante n'est pas
satisfaite, le type est disqualifié et son score divisé par quatre :

\begin{equation}
s_t^{\,\text{final}} =
\begin{cases}
s_t & \text{si toutes les règles bloquantes de } t \text{ sont satisfaites,}\\[6pt]
0{,}25 \times s_t & \text{sinon.}
\end{cases}
\label{eq:disqualification}
\end{equation}

Le choix de diviser plutôt que d'annuler est délibéré. Un type disqualifié reste visible
dans les alternatives avec son score pénalisé et le motif de sa disqualification. Une
opportunité de vision par ordinateur sans images annotées apparaît donc comme
« pertinente mais bloquée sur la vérité terrain », information bien plus utile qu'une
absence.

Le type retenu est celui de score final maximal ; en cas d'égalité, l'ordre alphabétique
du code technique départage, ce qui garantit un résultat reproductible. Les trois
suivants sont renvoyés comme alternatives.

\subsection{Enchaînement}

\begin{figure}[H]
\centering
\begin{tikzpicture}[node distance=0.55cm, scale=0.95, transform shape]
  \node[bloc, text width=4.6cm, fill=medicaSoft, draw=medicaGreen] (opp)
    {Opportunité\\\scriptsize 15 faits de qualification};
  \node[bloc, below=of opp, text width=4.6cm] (facts)
    {Extraction des faits\\\scriptsize surcharge par le questionnaire};
  \node[bloc, below=of facts, text width=6.2cm, draw=medicaTeal] (loop)
    {Pour chacun des 8 types :\\\scriptsize évaluer les règles, cumuler les poids satisfaits};
  \node[bloc, below=of loop, text width=6.2cm] (bloq)
    {Règle bloquante non satisfaite ?\\\scriptsize si oui : score $\times$ 0,25 et motif conservé};
  \node[bloc, below=of bloq, text width=6.2cm] (tri)
    {Tri décroissant des 8 scores};
  \node[blocvert, below=of tri, text width=7cm] (out)
    {\textbf{Recommandation}\\\scriptsize type retenu · confiance · justification rédigée\\
     \scriptsize règles satisfaites et manquées · 3 alternatives · garde-fous};
  \node[bloc, right=1.2cm of out, text width=3cm, draw=medicaAmber, fill=medicaBg, font=\scriptsize]
    (gxp) {Si \code{regulated\_gxp} : 3 garde-fous GxP ajoutés en tête};

  \draw[flechevert] (opp) -- (facts);
  \draw[flechevert] (facts) -- (loop);
  \draw[flechevert] (loop) -- (bloq);
  \draw[flechevert] (bloq) -- (tri);
  \draw[flechevert] (tri) -- (out);
  \draw[fleche, dashed, draw=medicaAmber] (gxp) -- (out);
\end{tikzpicture}
\caption{Enchaînement du moteur de qualification}
\label{fig:moteur-qualification}
\end{figure}

\subsection{Le cœur de l'algorithme}

\begin{lstlisting}[language=Python, caption={Évaluation d'un type de solution, extrait de services/recommendation.py}, label={lst:evaluate}]
def evaluate_type(ai_type, facts):
    rules = RULES[ai_type]
    total_weight = sum(r.weight for r in rules)
    fired_weight = 0.0
    disqualified_by = None

    for rule in rules:
        satisfied = bool(rule.predicate(facts))
        if satisfied:
            fired_weight += rule.weight
            fired.append(...)          # trace conservee pour l'explication
        else:
            missed.append(...)
            if rule.blocking and disqualified_by is None:
                disqualified_by = rule.label

    score = fired_weight / total_weight
    if disqualified_by:
        score *= 0.25                  # le type reste visible, fortement penalise
    return TypeEvaluation(ai_type, score, fired, missed, disqualified_by)
\end{lstlisting}

L'élément notable n'est pas le calcul, qui tient en dix lignes, mais le fait que les
listes \code{fired} et \code{missed} soient conservées. C'est ce qui permet à l'API de
répondre « voici pourquoi » et non seulement « voici quoi ». Ces listes sont persistées
en JSON dans \code{ai\_recommendation.fired\_rules}, ce qui rend la recommandation
reconstituable des mois plus tard, même si les règles ont changé entre-temps.

\subsection{Illustration}

Une opportunité d'inspection visuelle de blisters déclare des images disponibles et
annotées, une tâche répétitive, un volume mensuel de plusieurs milliers d'unités et un
périmètre GxP. Les cinq règles du type \emph{vision par ordinateur} sont satisfaites, le
score vaut 1,00. Le type \emph{IA classique} est disqualifié par la règle CLS-2, les
données n'étant pas structurées ; il figure en alternative avec un score divisé par
quatre et le motif « Données structurées (ERP, MES, LIMS) ». Le périmètre GxP ajoute en
tête des garde-fous la validation informatique, la traçabilité ALCOA+ et la revue
humaine obligatoire.

% =====================================================================
\section{L'algorithme de readiness (M5)}

\subsection{Formules}

Chaque dimension du diagnostic porte une note $n_i \in [0,5]$, un poids $p_i$ et une
confiance $c_i \in [0,1]$ déclarée par le répondant. Le score global est :

\begin{equation}
\text{readiness} = 100 \times
\frac{\displaystyle\sum_{i} n_i \, p_i \, c_i}{\displaystyle\sum_{i} 5 \, p_i \, c_i}
\label{eq:readiness}
\end{equation}

La confiance apparaît au numérateur \emph{et} au dénominateur. Une réponse mal étayée
pèse donc moins dans les deux termes : elle ne fait ni monter ni descendre
artificiellement le score, elle s'efface.

La moyenne d'un axe — technique ou métier — n'utilise que le poids :

\begin{equation}
\bar{n}_{\text{axe}} =
\frac{\displaystyle\sum_{i \in \text{axe}} n_i \, p_i}{\displaystyle\sum_{i \in \text{axe}} p_i}
\label{eq:moyenne-axe}
\end{equation}

La différence de traitement est intentionnelle. Le score global mesure une maturité
\emph{consolidée}, pour laquelle la solidité du constat compte. La moyenne d'axe décrit
le niveau constaté, indépendamment de la qualité de l'instruction.

\subsection{Dimensions non instruites}

Une confiance nulle signale une dimension non encore évaluée. Elle n'apporte ni au
numérateur ni au dénominateur — une dimension non mesurée ne vaut pas zéro — et elle
pèse uniquement sur le taux de complétion :

\begin{equation}
\text{complétion} = 100 \times
\frac{\text{dimensions instruites}}{\text{dimensions du référentiel}}
\end{equation}

Ce taux est toujours calculé, jamais saisi. La conséquence pratique est importante :
ajouter une dimension au référentiel fait baisser le taux de complétion de toutes les
campagnes, ce qui est exact, sans faire baisser mécaniquement leur score, ce qui serait
faux.

\subsection{Dimensions bloquantes et points de vigilance}

Une dimension instruite dont la note passe sous le seuil de la campagne — 2,5 par
défaut — est signalée. Sa qualification dépend de son caractère critique :

\begin{equation}
\text{bloquante} \iff (n_i < \theta) \wedge \text{critique}_i
\qquad
\text{vigilance} \iff (n_i < \theta) \wedge \neg\,\text{critique}_i
\end{equation}

Une dimension bloquante interdit le passage en production tant qu'elle n'est pas
remédiée. Un point de vigilance relève du plan de conduite du changement. Le niveau de
maturité affiché est $\lfloor \text{readiness} / 20 \rfloor$, plafonné à~5, sur une
échelle allant de « Inexistant » à « Optimisé ».

% =====================================================================
\section{La priorisation du portefeuille (M7)}

\subsection{Le score de priorité}

Six critères notés sur 100 sont combinés, puis une pénalité de risque est retranchée :

\begin{equation}
\text{score} = \sum_{k \in K} \omega_k \, v_k \; - \; \pi
\qquad\text{avec}\qquad
\sum_{k \in K} \omega_k = 1
\label{eq:priorite}
\end{equation}

Les pondérations par défaut sont données au tableau~\ref{tab:poids}.

\begin{table}[H]
\centering
\caption{Pondérations par défaut des six critères de priorisation}
\label{tab:poids}
\small
\begin{tabular}{@{}L{5cm} C{2cm} L{6.4cm}@{}}
\toprule
\textbf{Critère} & \textbf{Poids} & \textbf{Ce qu'il mesure} \\
\midrule
Impact métier             & 0,25 & Gain opérationnel attendu sur le processus \\
Impact financier / P\&L   & 0,20 & Effet sur le compte de résultat \\
Faisabilité technique     & 0,20 & Disponibilité des données et des compétences \\
Alignement stratégique    & 0,15 & Cohérence avec les priorités de l'organisation \\
Readiness organisationnel & 0,10 & Maturité du département porteur \\
Synergies                 & 0,10 & Réutilisation d'actifs partagés \\
\midrule
\textbf{Total}            & \textbf{1,00} & \\
\bottomrule
\end{tabular}
\end{table}

Ces poids sont ajustables depuis l'interface. Toute pondération fournie est renormalisée
à~1 avant usage, ce qui garantit que les scores restent comparables quelle que soit la
saisie. Un utilisateur qui porte l'impact financier de 0,20 à 0,60 sans toucher au reste
obtient bien une grille valide, dont les autres poids ont été proportionnellement
réduits.

\subsection{La pénalité de risque}

Le risque brut est le produit de la probabilité et de l'impact, tous deux notés de 1
à~5, soit une valeur de 1 à~25. Cette valeur est ramenée à l'échelle du score :

\begin{equation}
\pi = \frac{P \times I}{25} \times \Pi_{\max},
\qquad \Pi_{\max} = 15
\label{eq:penalite}
\end{equation}

Une opportunité au risque maximal perd donc 15~points sur 100. Le plafond est un choix
d'étalonnage : assez élevé pour déclasser une opportunité risquée, assez bas pour qu'un
risque maîtrisable n'annule pas un impact fort. Le score final est borné dans $[0,100]$.

\subsection{Rangs, vagues et quadrants}

Les opportunités sont triées par score décroissant, l'identifiant départageant les
égalités. Le rang est la position dans ce classement ; la vague est
$\lfloor (\text{rang}-1) / 6 \rfloor + 1$, la taille de vague étant paramétrable. Ce
découpage donne une lecture de séquencement plutôt qu'un simple palmarès : la vague~1
est ce que l'organisation lance maintenant, la vague~2 ce qu'elle prépare.

La matrice de portefeuille croise impact métier et faisabilité autour d'un seuil de~50,
définissant quatre quadrants : \emph{prioritaire} (fort impact, faisable),
\emph{quick win} (impact modéré, faisable), \emph{à réfléchir} (fort impact, difficile)
et \emph{à écarter}. La taille du point représente le bénéfice annuel attendu.

% =====================================================================
\section{Le pilotage de la valeur (M10)}

\subsection{Coût total de possession}

Les coûts sont séparés en non récurrents et récurrents. Sur un horizon de $H$ années :

\begin{equation}
\text{TCO} = \underbrace{\sum_{c \,\in\, \text{non récurrents}} m_c}_{\text{CAPEX}}
\; + \;
\underbrace{\Big(\sum_{c \,\in\, \text{récurrents}} m_c\Big)}_{\text{OPEX annuel}} \times H
\label{eq:tco}
\end{equation}

\subsection{Gains et retour sur investissement}

Chaque bénéfice porte un montant annuel et un taux de réalisation, qui traduit la
prudence de l'hypothèse. Le gain annuel retenu est pondéré par ce taux :

\begin{equation}
G = \sum_{b} m_b \, \rho_b
\qquad\text{puis}\qquad
\text{ROI}\,(\%) = \frac{G \times H - \text{TCO}}{\text{TCO}} \times 100
\label{eq:roi}
\end{equation}

\subsection{Délai de retour}

\begin{equation}
\text{payback (mois)} = \frac{\text{CAPEX}}{\big(G - \text{OPEX annuel}\big) / 12}
\label{eq:payback}
\end{equation}

Le numérateur est le seul CAPEX, pas le TCO. Le payback répond à la question « en
combien de mois l'investissement initial est-il récupéré par le gain annuel net
d'exploitation ? ». Y faire figurer trois années d'OPEX qui ne sont pas encore engagées
allongerait artificiellement le délai. Le calcul n'est produit que si le gain net annuel
et le CAPEX sont strictement positifs ; sinon le délai est déclaré non calculable plutôt
qu'affiché à zéro.

\subsection{Valeur promise et valeur réalisée}

Les bénéfices sont rattachés à des indicateurs, dont les mesures successives sont
horodatées. Le taux d'atteinte d'un indicateur est
$(\text{courant} - \text{référence}) / (\text{cible} - \text{référence}) \times 100$,
borné à 150~\% pour absorber les dépassements sans fausser les agrégats. C'est cette
comparaison qui ferme la boucle du parcours : le business case cesse d'être un document
d'obtention de budget pour devenir un engagement suivi.

% =====================================================================
\section{Sécurité et isolation}

\subsection{Chaîne de contrôle d'une requête}

\begin{figure}[H]
\centering
\begin{tikzpicture}[node distance=0.5cm, scale=0.95, transform shape]
  \node[bloc, text width=6cm, fill=medicaSoft, draw=medicaGreen] (req)
    {Requête HTTP\\\scriptsize en-tête \code{Authorization: Bearer}};
  \node[bloc, below=of req, text width=6cm, draw=medicaTeal] (jwt)
    {Vérification du jeton\\\scriptsize signature HS256, expiration};
  \node[bloc, below=of jwt, text width=6cm, draw=medicaTeal] (user)
    {Chargement de l'utilisateur\\\scriptsize compte actif ? sinon 401};
  \node[bloc, below=of user, text width=6cm, draw=medicaTeal] (perm)
    {Permission requise par l'endpoint\\\scriptsize absente $\rightarrow$ \textbf{403}};
  \node[bloc, below=of perm, text width=6cm, draw=medicaTeal] (tenant)
    {Filtre \code{organization\_id} issu du jeton\\\scriptsize jamais de la requête};
  \node[bloc, below=of tenant, text width=6cm, draw=medicaTeal] (scope)
    {Périmètre de l'utilisateur\\\scriptsize lecture hors périmètre $\rightarrow$ \textbf{404}\\
     \scriptsize écriture hors périmètre $\rightarrow$ \textbf{403}};
  \node[blocvert, below=of scope, text width=6cm] (exec)
    {Exécution et journalisation\\\scriptsize entrée \code{audit\_log} si écriture};

  \foreach \a/\b in {req/jwt, jwt/user, user/perm, perm/tenant, tenant/scope, scope/exec}
    \draw[flechevert] (\a) -- (\b);

  \node[bloc, right=1.2cm of perm, text width=3.1cm, fill=medicaBg, draw=medicaRed, font=\scriptsize]
    (refus) {Refus\\ 401 · 403 · 404\\ aucune donnée renvoyée};
  \draw[fleche, draw=medicaRed] (perm.east) -- (refus.west);
  \draw[fleche, draw=medicaRed] (scope.east) -| (refus.south);
  \draw[fleche, draw=medicaRed] (jwt.east) -| (refus.north);
\end{tikzpicture}
\caption{Contrôles traversés par toute requête authentifiée}
\label{fig:securite}
\end{figure}

\subsection{Authentification}

Les mots de passe sont hachés avec bcrypt ; aucun mot de passe en clair n'est stocké ni
journalisé. La connexion émet un jeton JWT signé en HS256, dont la charge utile contient
l'identifiant de l'utilisateur, celui de son organisation, son rôle principal, la date
d'émission et la date d'expiration. La durée de validité par défaut est de 720~minutes.

Deux points de configuration méritent attention avant une mise en production : la clé de
signature porte une valeur par défaut explicitement nommée
\code{dev-secret-key-change-me}, et le compte superutilisateur initial a un mot de passe
par défaut. Les deux doivent être redéfinis par variable d'environnement.

\subsection{Autorisation}

L'autorisation repose sur deux dépendances FastAPI centrales.
\code{require\_permission} vérifie qu'un utilisateur détient une permission et lève 403
sinon. \code{get\_scoped\_or\_404} charge une entité en appliquant simultanément le
filtre de tenant et le filtre de périmètre.

Le choix des codes de retour est délibéré. Une lecture hors périmètre répond~404, non
403 : répondre 403 confirmerait l'existence de la ressource, ce qui constitue déjà une
fuite d'information. Une écriture hors périmètre répond 403, car l'utilisateur voit
l'objet mais n'a pas le droit d'agir dessus.

\begin{remarque}[Le frontend ne protège rien]
Le masquage d'un bouton dans l'interface n'est pas un contrôle d'accès. Il améliore
l'ergonomie en n'affichant que les actions possibles. La protection réelle est
exclusivement côté serveur. Les tests d'autorisation appellent les endpoints
directement, sans passer par l'interface, précisément pour vérifier ce point.
\end{remarque}

\subsection{Isolation entre organisations}

L'identifiant d'organisation est lu dans le jeton, jamais dans le corps ou les
paramètres de la requête. Un utilisateur ne peut donc pas se déclarer membre d'une autre
organisation, même en forgeant sa requête, sans forger le jeton, ce qui exigerait la clé
de signature.

La permission \code{platform.admin} est la seule à franchir cette frontière. Elle ne
peut pas être placée par un administrateur d'entreprise dans un rôle qu'il crée : une
vérification explicite l'en empêche.

\subsection{Traçabilité}

Toute écriture significative produit une entrée dans \code{audit\_log}, conservant
l'auteur, son adresse électronique au moment de l'action, l'horodatage, le type et
l'identifiant de l'entité, un libellé lisible, l'ancienne et la nouvelle valeur en JSON,
et l'adresse IP. L'API n'expose aucune mise à jour ni suppression sur cette table : le
journal est en insertion seule. L'adresse électronique est dupliquée dans l'entrée
précisément pour que la trace reste lisible si le compte est supprimé par la suite.
